\documentclass[a4paper,12pt]{article}
\usepackage[utf8]{inputenc}
\usepackage[ngerman]{babel}
\usepackage{amsmath,amssymb}
\usepackage{geometry}
\geometry{margin=2.5cm}

\title{Aufgabe 8 -- Parametrisches Gleichungssystem}
\author{}
\date{}

\begin{document}
\maketitle

\section*{Gleichungssystem}

\[
\begin{cases}
(k+1)x + (k^2-2k-3)y = k^2-k+2 \\
(k-1)x + (k^2-2k+1)y = k^2-k
\end{cases}
\]

\textbf{Hinweis:} Es werden \textbf{keine Determinanten} verwendet, sondern nur Gauß-Elimination mit dem Parameter $k$.

\section*{Faktorisierung der Koeffizienten}

Zunächst faktorisieren wir, um die Struktur besser zu erkennen:
\begin{align*}
k^2-2k-3 &= (k+1)(k-3) \\
k^2-2k+1 &= (k-1)^2
\end{align*}

\section*{Erweiterte Koeffizientenmatrix}

\[
\left(\begin{array}{cc|c}
k+1 & (k+1)(k-3) & k^2-k+2 \\
k-1 & (k-1)^2 & k^2-k
\end{array}\right)
\]

\section*{Gauß-Elimination}

Wir wollen die erste Spalte in Zeile 2 eliminieren. Um ganzzahlig zu bleiben, rechnen wir:

\textbf{Schritt 1:} $Z_2 \leftarrow (k+1) \cdot Z_2 - (k-1) \cdot Z_1$

\textbf{Neue Zeile 2, Spalte für $y$:}
\begin{align*}
(k+1)(k-1)^2 - (k-1)(k+1)(k-3) &= (k+1)(k-1)\big[(k-1) - (k-3)\big] \\
&= (k+1)(k-1) \cdot 2 \\
&= 2(k^2-1)
\end{align*}

\textbf{Neue Zeile 2, rechte Seite:}
\begin{align*}
(k+1)(k^2-k) - (k-1)(k^2-k+2) &= k^3 - k^2 + k^2 - k - (k^3 - k^2 + 2k - k^2 + k - 2) \\
&= k^3 - k - k^3 + 2k^2 - 3k + 2 \\
&= 2k^2 - 4k + 2 \\
&= 2(k-1)^2
\end{align*}

Die Matrix nach der Elimination:

\[
\left(\begin{array}{cc|c}
k+1 & (k+1)(k-3) & k^2-k+2 \\
0 & 2(k^2-1) & 2(k-1)^2
\end{array}\right)
\]

Da $k^2-1 = (k-1)(k+1)$, lautet Zeile 2:

\[
2(k-1)(k+1) \cdot y = 2(k-1)^2
\]

\section*{Fallunterscheidung}

\subsection*{Fall 1: $k \neq 1$ und $k \neq -1$}

Wir können durch $2(k-1)$ kürzen:

\[
(k+1) \cdot y = (k-1) \quad \Rightarrow \quad y = \frac{k-1}{k+1}
\]

Aus Zeile 1 (mit $k+1 \neq 0$):
\begin{align*}
(k+1)x &= k^2-k+2 - (k+1)(k-3) \cdot \frac{k-1}{k+1} \\
&= k^2-k+2 - (k-3)(k-1) \\
&= k^2-k+2 - (k^2-4k+3) \\
&= 3k - 1
\end{align*}

\[
x = \frac{3k-1}{k+1}
\]

\[
\boxed{\text{Für } k \neq \pm 1: \text{ eindeutige Lösung } \left(\frac{3k-1}{k+1},\; \frac{k-1}{k+1}\right)}
\]

\subsection*{Fall 2: $k = 1$}

Setze $k = 1$ direkt in das Gleichungssystem ein:
\begin{align*}
2x + (-4)y &= 2 \\
0 \cdot x + 0 \cdot y &= 0
\end{align*}

Zeile 2 wird zu $0 = 0$ (keine Information). Aus Zeile 1:
\[
2x - 4y = 2 \quad \Rightarrow \quad x = 1 + 2y
\]

Die Variable $y$ ist frei. Setze $y = s$ mit $s \in \mathbb{R}$.

\[
\boxed{\text{Für } k = 1: \text{ unendlich viele Lösungen } (x, y) = (1 + 2s,\; s), \quad s \in \mathbb{R}}
\]

\subsection*{Fall 3: $k = -1$}

Setze $k = -1$ direkt in das Gleichungssystem ein:
\begin{align*}
0 \cdot x + (1+2-3)y &= 1+1+2 \\
(-2)x + (1+2+1)y &= 1+1
\end{align*}

Also:
\begin{align*}
0 \cdot x + 0 \cdot y &= 4 \\
-2x + 4y &= 2
\end{align*}

Zeile 1 ergibt $0 = 4$ -- ein \textbf{Widerspruch}.

\[
\boxed{\text{Für } k = -1: \text{ keine Lösung (inkonsistent)}}
\]

\section*{Zusammenfassung}

\[
\begin{cases}
k \neq \pm 1: & \text{eindeutige Lösung} \\
k = 1: & \text{unendlich viele Lösungen} \\
k = -1: & \text{keine Lösung}
\end{cases}
\]

\end{document}
